% =====================================================================
%  CHAPTER 10: 学校健康教育（对应大纲第十章）
%  内容来源：往届笔记第十五章，参考PPT第十七章
% =====================================================================
\chapter{学校健康教育}
\setcounter{chapter}{10}
\chapterintro{
\textbf{熟悉：}学校健康教育的概念、基本内容、专题教育、学校性教育。学校健康教育的实施、主要理论与应用、评价类型、评价方法。
}

% ===== 熟悉内容 =====
\section{概述}

\subsection{基本概念}

\begin{defbox}
\textbf{学校健康教育（school health education）：}以促进学生健康为核心的教育活动与过程。通过有计划、有组织、多种形式的教育教学活动，使学生掌握卫生保健知识，增强学生自我保健意识，养成科学、文明、健康的生活方式和行为习惯，达到预防疾病、增进身心健康、全面提高学生健康水平的目的。

\textbf{健康促进学校（health promoting schools）：}学校内所有成员为维护和促进师生健康共同努力，制定促进师生健康的规章制度，提供完整、积极的经验和知识结构，包括设置正式和非正式健康课程，创造安全健康的学校环境，提供适宜的健康服务。

\textbf{健康素养（health literacy）：}个人获取和正确理解健康信息与服务，并运用这些信息和服务作出正确判断，维护和促进自身健康的能力。

\textbf{学生健康素养：}学生获取和正确理解健康信息与服务，并运用这些信息和服务作出正确判断，维护和促进自身健康的能力。
\end{defbox}

\subsection{基本内容}

\begin{keybox}
\textbf{5个领域：}

\begin{enumerate}[leftmargin=*]
    \item \imp{健康行为与生活方式}
    \item \imp{疾病预防}
    \item \imp{心理健康}
    \item \imp{生长发育与青春期保健}
    \item \imp{安全应急与避险}
\end{enumerate}

\textbf{5个水平：}

\begin{enumerate}[leftmargin=*]
    \item 水平一（小学1-2年级）
    \item 水平二（小学3-4年级）
    \item 水平三（小学5-6年级）
    \item 水平四（初中7-9年级）
    \item 水平五（高中10-12年级）
\end{enumerate}
\end{keybox}

\subsection{专题教育}

\begin{keybox}
\textbf{学校健康教育三大专题教育：}

\begin{enumerate}[leftmargin=*]
    \item \imp{预防艾滋病专题教育}
    \item \imp{毒品预防专题教育}
    \item \imp{预防烟草教育}
\end{enumerate}
\end{keybox}

\section{学校健康教育与健康促进实施}

\subsection{教学方法}

\begin{keybox}
\textbf{1. 直接式}

（1）\textbf{传统式：}课堂讲授；讲座；示教。

（2）\textbf{参与式：}小组讨论和案例分析；头脑风暴；角色扮演、小品和游戏活动；辩论和演讲；同伴教育。

\textbf{2. 间接式：}大众媒体；视听手段；网络系统性学习。
\end{keybox}

\subsection{学校健康教育组织实施}

\begin{tipbox}
\textbf{组织实施要点：}

\begin{enumerate}[leftmargin=*]
    \item \textbf{课堂教学}——将健康教育纳入课程体系，保证课时。
    \item \textbf{教师培训}——提高教师健康教育能力和素养。
    \item \textbf{学校支持性环境}——创建有利于健康的学校物质和社会环境。
\end{enumerate}
\end{tipbox}

\section{健康教育与健康促进主要理论与学校应用}

\subsection{生活技能理论}

\begin{defbox}
\textbf{生活技能（life skills）：}一个人的心理社会能力。

\textbf{生活技能教育（life skills-based education）：}促进心理社会能力在适宜的文化背景下实践并得到发展，促进个体和社会发展，预防可能出现的健康和社会问题，保障人权。
\end{defbox}

\begin{keybox}
\textbf{生活技能核心能力（5对）：}

\begin{enumerate}[leftmargin=*]
    \item \imp{自我认识能力——同理能力}
    \item \imp{有效交流能力——人际关系能力}
    \item \imp{调节情绪能力——缓解压力能力}
    \item \imp{创造性思维能力——批判性思维能力}
    \item \imp{决策能力——解决问题能力}
\end{enumerate}
\end{keybox}

\subsection{组织改变理论}

\begin{keybox}
\textbf{组织改变理论模型（theory of organizational change）：}通过人群所在组织的改变、规章制度和管理模式的完善，实现对组织内全体人群的干预。

\textbf{4个阶段：}

\begin{enumerate}[leftmargin=*]
    \item \textbf{确立问题或认知阶段：}对问题的认知与分析，寻求解决问题的方法和评价。
    \item \textbf{初期行动阶段：}制订项目规划和实施计划，为开始改变准备。
    \item \textbf{执行阶段：}创新的执行阶段。
    \item \textbf{制度化阶段：}新政策等变为确定的、新的目标价值融入组织里。
\end{enumerate}
\end{keybox}

\subsection{创新扩散理论}

\begin{keybox}
\textbf{创新扩散（diffusion of innovation, DI）：}一项新事物通过一定传播渠道在整个社区或某个人群内扩散，逐渐为社区成员或该人群成员了解与采用的过程。

\textbf{5个阶段：}

\begin{enumerate}[leftmargin=*]
    \item \textbf{认知阶段：}个体首次接触新事物，但是缺少相关信息。
    \item \textbf{说服阶段：}个体对新事物产生兴趣，积极寻求相关信息。
    \item \textbf{决策阶段：}评估使用新事物的优势和劣势，决定采用或拒绝。
    \item \textbf{实施阶段：}个体不同程度采用了新事物。
    \item \textbf{确认阶段：}个人定型，决定继续使用创新。
\end{enumerate}

\textbf{5种采纳者类型：}

\begin{enumerate}[leftmargin=*]
    \item \textbf{创新者：}人群中最先接受信息。
    \item \textbf{早期采用者：}较易接受新观念、态度慎重、常具领导力。
    \item \textbf{中期采用者：}慎重、深思熟虑。
    \item \textbf{后期采用者：}持怀疑态度，等多数人接受后才采用。
    \item \textbf{迟缓者：}观念保守、坚持习惯，不到万不得已不愿接受创新。
\end{enumerate}
\end{keybox}

\subsection{其他理论}

\begin{tipbox}
\textbf{理性行动理论（theory of reasoned action, TRA）/ 理性行为理论：}分析态度如何有意识影响个体行为的理论。

\textbf{大众意见领袖理论：}
\begin{itemize}
    \item \textbf{大众意见领袖（popular opinion leaders, POLs）/ 舆论领袖：}大众信息传播中信息从源头到一般受众的中间处理环节或节点。
    \item \textbf{信息传播模式：}大众传播→意见领袖→一般受众。
\end{itemize}
\end{tipbox}

\section{学校健康教育与健康促进评价}

\subsection{评价类型}

\begin{keybox}
\textbf{1. 形成评价}

\textbf{2. 过程评价}
\begin{enumerate}[leftmargin=*]
    \item 健康教育活动评价
    \item 学校卫生服务评价
    \item 学校环境评价
\end{enumerate}

\textbf{3. 效果评价}
\begin{enumerate}[leftmargin=*]
    \item 健康教育活动效果评价
    \item 卫生服务效果评价
\end{enumerate}
\end{keybox}

\subsection{评价方法与维度}

\begin{keybox}
\textbf{评价方法：}

\begin{enumerate}[leftmargin=*]
    \item \textbf{观察法}
    \item \textbf{个人访谈、小组访谈和资料与案例分析}
    \item \textbf{调查问卷：}设置封闭式和开放式问题，可用于知识、态度和行为的评价。
\end{enumerate}

\textbf{评价维度与指标：}知识、态度、技能、行为。
\end{keybox}

% ----- 复习题 -----
\reviewcontent{
\section{复习题}

\begin{enumerate}[leftmargin=*, itemsep=8pt]
    \item \textbf{【名词解释】}学校健康教育；健康促进学校；健康素养；生活技能；创新扩散；大众意见领袖
    \item \textbf{【简答题】}简述学校健康教育的基本内容（5个领域和5个水平）。
    \item \textbf{【简答题】}简述学校健康教育的三大专题教育。
    \item \textbf{【简答题】}简述生活技能理论的5对核心能力。
    \item \textbf{【简答题】}简述组织改变理论的4个阶段。
    \item \textbf{【简答题】}简述创新扩散理论的5个阶段和5种采纳者类型。
    \item \textbf{【简答题】}简述学校健康教育与健康促进的评价类型和方法。
\end{enumerate}

\subsection*{参考答案要点}

\begin{enumerate}[leftmargin=*, itemsep=5pt]
    \item \textbf{学校健康教育}：以促进学生健康为核心的教育活动与过程，通过有计划、有组织的教育教学活动，使学生掌握卫生保健知识，养成科学生活方式。\textbf{健康促进学校}：学校内所有成员为维护和促进师生健康共同努力，制定规章制度，提供完整积极的经验和知识结构。\textbf{健康素养}：个人获取和正确理解健康信息与服务，并运用这些信息和服务作出正确判断，维护和促进自身健康的能力。\textbf{生活技能}：一个人的心理社会能力。\textbf{创新扩散}：一项新事物通过一定传播渠道在社区或人群中扩散，逐渐被了解与采用的过程。\textbf{大众意见领袖}：大众信息传播中信息从源头到一般受众的中间处理环节或节点。

    \item 5个领域：(1)健康行为与生活方式；(2)疾病预防；(3)心理健康；(4)生长发育与青春期保健；(5)安全应急与避险。5个水平：(1)水平一（小学1-2年级）；(2)水平二（小学3-4年级）；(3)水平三（小学5-6年级）；(4)水平四（初中7-9年级）；(5)水平五（高中10-12年级）。

    \item (1)预防艾滋病专题教育；(2)毒品预防专题教育；(3)预防烟草教育。

    \item (1)自我认识能力——同理能力；(2)有效交流能力——人际关系能力；(3)调节情绪能力——缓解压力能力；(4)创造性思维能力——批判性思维能力；(5)决策能力——解决问题能力。

    \item (1)确立问题或认知阶段：对问题的认知与分析；(2)初期行动阶段：制订项目规划和实施计划；(3)执行阶段：创新的执行；(4)制度化阶段：新政策变为确定的目标价值融入组织。

    \item 5个阶段：(1)认知→(2)说服→(3)决策→(4)实施→(5)确认。5种类型：(1)创新者→(2)早期采用者→(3)中期采用者→(4)后期采用者→(5)迟缓者。

    \item 评价类型：(1)形成评价；(2)过程评价（健康教育活动/卫生服务/学校环境评价）；(3)效果评价（健康教育活动/卫生服务效果评价）。评价方法：观察法、访谈和案例分析、调查问卷。评价维度：知识、态度、技能、行为。
\end{enumerate}
}

\newpage
